% from ega3-6-fr.tex
\oldpage[III]{137}

\begin{env}[6.1.1]
Soient $f:X\to Y$ un morphisme de préschémas, $\mathscr F$ un
$\mathscr O_X$-Module quasi-cohérent. Dans l'étude des « images directes
supérieures » $\mathrm R^n f_*(\mathscr F)$, on est amené à considérer le
problème général suivant : étant donné un morphisme « changement de base »
$g:Y'\to Y$, on pose
\[
  X'=X_{(Y')}=X\times_Y Y',\qquad
  \mathscr F'=\mathscr F\otimes_Y\mathscr O_{Y'},\qquad
  f'=f_{(Y')}:X'\to Y',
\]
et l'on se propose d'avoir des renseignements sur les images directes
supérieures $\mathrm R^n f'_*(\mathscr F')$ (en supposant connus les
$\mathrm R^n f_*(\mathscr F)$). On voit aisément (cf. par exemple (7.7.2))
que l'on est ramené à étudier les variations de
$\mathrm R^n f_*(\mathscr F\otimes_{\mathscr O_Y}\mathscr G)$ pour un
$\mathscr O_Y$-Module quasi-cohérent variable $\mathscr G$, autrement dit le
foncteur
\[
  \mathscr G\longmapsto
  \mathrm R^n f_*(\mathscr F\otimes_{\mathscr O_Y}\mathscr G).
\]
Si $\mathscr F$ est plat sur $Y$, le foncteur
$\mathscr G\mapsto\mathscr F\otimes_{\mathscr O_Y}\mathscr G$ est exact
(0\textsubscript{I}, 6.7.4) et par suite le foncteur composé
$\mathscr G\mapsto\mathrm R^n f_*(\mathscr F\otimes_{\mathscr O_Y}\mathscr G)$
est encore un foncteur cohomologique. Mais il n'en est plus de même dans le
cas général; pour pouvoir appliquer les méthodes cohomologiques, on est
conduit à substituer à
$\mathscr G\mapsto\mathrm R^n f_*(\mathscr F\otimes_{\mathscr O_Y}\mathscr G)$
d'autres foncteurs, qui cette fois sont toujours des foncteurs
cohomologiques. Ces foncteurs, qui généralisent les foncteurs « Tor » de la
théorie des modules, sont définis aux n\textsuperscript{os} 6.3 à 6.7; il y a
d'ailleurs deux telles généralisations, l'une « locale » et l'autre
« globale », reliées par des suites spectrales qui seront discutées au
n\textsuperscript{o} 6.7; comme application de ces suites spectrales, on
obtient en particulier, sous certaines conditions, une « formule de Künneth »
exprimant
\[
  \mathrm R^n(f_1\times f_2)_*(\mathscr F_1\otimes_Y\mathscr F_2)
\]
à l'aide des images directes supérieures
$\mathrm R^p f_{1*}(\mathscr F_1)$ et $\mathrm R^q f_{2*}(\mathscr F_2)$.
D'autres suites spectrales (6.8) généralisent les suites spectrales
d'associativité du foncteur « Tor » de modules; enfin, le problème du
changement de base conduit lui aussi à des suites spectrales (6.9).
\end{env}

\begin{env}[6.1.2]
En outre, on constate (6.10) que les foncteurs cohomologiques
$\mathscr G\mapsto\mathscr T_\bullet(\mathscr G)$ ainsi définis sont
localement (sur $Y$) du type
\[
  \mathscr G\longmapsto
  \mathscr H_\bullet(\mathscr L_\bullet\otimes_Y\mathscr G),
\]
où $\mathscr L_\bullet$ est un complexe de $\mathscr O_Y$-modules localement
libres (défini à une homotopie près) et $\mathscr H_\bullet$ l'homologie. Il y
a alors intérêt à oublier la situation particulière qui a donné naissance à
$\mathscr T_\bullet$, et à étudier de

